\begin{textsample}{POS Dim 5 – global\_south\_2024\_12 – Score 15.00 – global\_south\_2024\_12/119053052.txt}  \label{ex:f5_pos_050}
Israel denies HRW ' s \textbf{genocide} claims in Gaza A New York-based rights group, Human Rights Watch ( HRW ) has accused Israel of \textbf{committing}"\textbf{acts} of \textbf{genocide}"in the Gaza Strip by damaging water infrastructure and cutting off supplies to civilians.
Israel has since denied the accusations as"appalling lies."In a new report which focused specifically on water, HRW detailed that what it said were deliberate efforts by Israeli authorities"of a systematic nature"to deprive Gazans of water, which had"likely caused thousands of deaths and will likely continue to cause deaths."Recall that the war in Gaza was sparked by Hamas ' s October 7, 2023 attack on Israel that resulted in the deaths of 1,208 people, mostly civilians.
Since then, Israel ' s retaliatory offensive has killed at least 45,097 people in Gaza, a majority of them civilians, according to figures from the Hamas-run territory ' s health ministry that the United Nations considers reliable."Since October 2023, Israeli authorities have deliberately obstructed Palestinians ' access to the adequate amount of water required for survival in the Gaza Strip,"the report said."Human Rights Watch is once more spreading its blood libels to promote its anti-Israel propaganda.
This report is full of lies that are appalling even when compared to HRW ' s already low standards."Since the beginning of the war, Israel has facilitated the \textbf{continuous} flow of water and humanitarian aid into Gaza, despite operating under constant attacks of Hamas terror organisation,"the Foreign Ministry said in a statement.
The HRW report detailed what the group said was the intentional damage to water and sanitation infrastructure, including solar panels powering treatment plants, a reservoir and a spare parts warehouse, as well as the blocking of fuel for generators.
The report concluded that"Israeli authorities \textbf{intentionally} inflicted on the Palestinian population in Gaza ' conditions of life calculated to bring about its physical destruction in whole or in part. '"This, it said, amounted to the war \textbf{crime} of"extermination"and to"\textbf{acts} of \textbf{genocide}."Israel ' s foreign ministry had noted that it has"ensured water infrastructure, including the continued operation of four water pipelines and water pumping and desalination facilities, which remain operational."HRW stopped short of saying Israel was \textbf{committing} outright"\textbf{genocide}."Under \textbf{international} law, proving \textbf{genocide} requires evidence of specific \textbf{intent}, which experts say is very difficult.
Speaking at a briefing on the report, Lama Faqih, director of HRW ' s Middle East and North Africa division, said that in the absence of"a clear articulated plan"to \textbf{commit} \textbf{genocide}, the \textbf{International} \textbf{Court} of Justice ( ICJ ) might find that the evidence meets the"very strict threshold"of \textbf{reasonable} inference of \textbf{genocidal} \textbf{intent}.
HRW pointed to a statement by then-defence minister Yoav Gallant in October 2023, when he declared a"complete siege"and said :"No electricity, no food, no water, no gas -- it ' s all closed."Israel is facing a \textbf{case} brought by South Africa at the \textbf{International} \textbf{Court} of Justice last December, arguing that the war in Gaza breached the 1948 United Nations \textbf{Genocide} \textbf{Convention}, an accusation Israel has strongly denied.
On December 5, \textbf{Amnesty} \textbf{International} accused Israel of \textbf{committing} \textbf{genocide} in the Gaza Strip, drawing a furious reaction from the government.
The HRW report, drawn up over nearly a year, is based on interviews with dozens of Gazans, staff at water and sanitation facilities, medics and aid workers, as well as satellite imagery, photographs, videos and data \textbf{analysis}.
It said Israeli authorities did not reply to requests for information.
The lack of water left Gazans vulnerable to water-borne diseases and complications, such as infected wounds and the inability to heal due to dehydration, HRW said.

% matched lemmas: act, amnesty, analysis, case, commit, continuous, convention, court, crime, genocidal, genocide, intent, intentionally, international, reasonable
\end{textsample}
